%%%%%%%%%%%%%%%%%%%%%%%%%%%%%%%%%%%%%%%%%%%%%%%%%%%%%%%%%%%%%%%%%%%%%%
% How to use writeLaTeX: 
%
% You edit the source code here on the left, and the preview on the
% right shows you the result within a few seconds.
%
% Bookmark this page and share the URL with your co-authors. They can
% edit at the same time!
%
% You can upload figures, bibliographies, custom classes and
% styles using the files menu.
%
%%%%%%%%%%%%%%%%%%%%%%%%%%%%%%%%%%%%%%%%%%%%%%%%%%%%%%%%%%%%%%%%%%%%%%

\documentclass[12pt]{article}

\usepackage{sbc-template}

\usepackage{graphicx,url}

%\usepackage[brazil]{babel}   
\usepackage[utf8]{inputenc}  

     
\sloppy

\title{Título do relatório de estágio que identifique a atividade realizada}

\author{Nome do Aluno\inst{1}, Nome do Orientador\inst{1} }


\address{Universidade Federal de Goiás - Regional Jataí - UFG REJ\\
  BR 364, km 195, nº 3800 CEP 75801-615 - Jataí - Goiás
  \email{e-mail.do.aluno@gmail.com.br, e-mail.do.orientador@gmail.com}
}

\begin{document} 

\maketitle

\begin{abstract}
  This report describes the activities...
\end{abstract}

\begin{resumo} 
  Este relatório descreve as atividades realizadas no estágio supervisionado obrigatório, item obrigatório para a conclusão da disciplina de Estágio do Curso de Ciências da Computação da Universidade Federal de Goiás, Regional Jataí (UFG-REJ). Atividades estas que foram realizadas na empresa ... que tem como principal atividade econômica ...e as atividades internas que utilizam recursos computacionais é a ... as quais foram observadas com o intuíto de apresentar melhorias e inovações com foco em .... 
  
\end{resumo}


\section{Introdução}\label{sec:introducao}

No presente documento serão apresentados e descritos as atividades desenvolvidas durante o estágio de 128 horas, realizado na empresa identificada do local do estágio, localizada na cidade a ser identificada, no período registrado conforme as folhas de registro de frequência (aqui pode colocar o período em data), sendo a atividade de estágio na modalidade de observação. O estágio foi focado na Análise e Proposta de Intervenção nas atividades relacionadas a recursos tecnológicos, (descrever as que foram utilizadas) como por exemplo: desenvolvimento de sistemas, configuração de redes de computadores, serviços de gerenciamento de banco de dados, renovação e atualização de hardwares e demais inovações computacionais relacionadas com possíveis contribuições para aperfeiçoar as atividades da referida empresa.
	
    O estágio foi orientado por ..., doutor/doutora/mestre em Ciência da Computação pela ... e supervisionado por ..., graduado em (obrigatório..) pela ... que é sócio proprietário da empresa ou gerente do setor ou responsável pelo setor ou administrador da empresa onde foi realizado o estágio.
    
      O referido relatório, está organizado em seções, as quais, serão apresentadas da seguinte forma: seção 2, um relato da importância do estágio, no âmbito pessoal como profissional, e ainda, uma breve análise da empresa. Na seção 3, embasamento teórico do relatório, descrevendo as teorias de base utilizadas para propor a intervenção.  Na seção 4, descrição mais detalhada sobre a observação realizada e a identificação da intervenção a ser proposta. Na seção 5, descrição e detalhamento da intervenção proposta. Na seção 6 as considerações finais sobre as atividades realizadas no estágio supervisionado e sobre as perspectivas futuras da proposta de intervenção e por último, a seção 7 as referências bibliográficas.

\section{O Estágio supervisionado e o local de atuação} \label{sec:oestagio}
Nesta seção será apresentado um breve estudo sobre o quê é o estágio supervisionado obrigatório, em qual curso o mesmo está sendo realizado, terminando o parágrafo apresentando a empresa do local do estágio.

\subsection{Histórico da empresa}
Nesta subseção deve descrever qual é o seguimento da empresa, qual é a importância do seguimento para a região de atuação do aluno, como se apresenta este seguimento frente a futura profissão do estagiário.

	Após situar o leitor sobre o seguimento de atuação, descrever de forma mais específica o histórico da empresa e quais são suas atividades principais e por fim qual setor da empresa o estagiário fez a sua observação, descrevendo de forma detalhada como é realizado as tarefas no referido setor, terminando a seção com uma sugestão de melhoria no setor, que será o foco principal do relatório nas próximas seções. Podendo descrever quantos funcionários a empresa possui, quantos computadores e/ou quais recursos computacionais utilizados (softwares em geral), dando atenção a itens que irão fazer parte da proposta de intervenção.

\section{Embasamento Teórico} \label{sec:embasteorico}

Nesta seção, será apresentado todo o embasamento teórico necessário para o desenvolvimento das atividades. Pode-se destacar os métodos utilizados para Análise e Desenvolvimento de Sistemas, as técnicas a serem utilizadas para pesquisa e avaliação de procedimentos e demais recursos que forem utilizados na proposta.


\subsection{Linguagem de Programação Java}

Lançado em 1996, o Java é uma ótima opção para se programar um sistema, pois oferece várias vantagens. Uma delas, é seu alto nível de segurança, por ser uma linguagem de programação usada em ambientes de rede, onde encontram-se variados sistemas operacionais e arquiteturas de CPU, por exemplo.

\begin{quote}
Segundo Horstmann e Cornell (2010, p.3) "Desde o início, o Java foi projetado para tornar certos tipos de ataques impossíveis, entre eles: estouro da pilha em tempo de execução - um ataque comum de vermes e vírus -, corrupção de memória fora do seu próprio espaço de processo e leitura e gravação de arquivos sem permissão".
\end{quote}

Conforme o passar dos anos, outros recursos foram adicionados para enriquecer ainda mais a segurança do Java.

	Outra vantagem da utilização dessa linguagem é o fato dela ser neutra em relação à arquitetura, ou seja, seu software pode ser compilado em qualquer máquina, por mais que esta não possua uma arquitetura tão boa. Tudo isso é graças ao compilador Java que gera instruções bytecodes, que são projetadas para serem de fácil interpretação em qualquer máquina. Entretando, esse processo pode ser tornar mais lento que executar instruções de máquina em uma velocidade total. Mas as máquinas virtuais possuem um processo chamado compilação just-in-time, que transforma as sequências de bytecodes mais usadas em código de máquina. Esse processo fez tanto sucesso que a própria plataforma .NET da Microsoft adquiriu uma máquina virtual.
    
	Assim como outras linguagens de programação, o Java também possui Programação Orientada a Objetos (POO), onde, na utilização deste, o desenvolvedor preocupa-se mais com os dados (objetos), do que com o código fonte. Esse modo de programação é muito dominante hoje em dia, substituindo a programação procedimental "estruturada".
    
    Netbeans...

\section{Descrição das Atividades Realizadas}
Dentre as várias atividades da empresa a atividade de ... foi escolhida como objeto de estudo deste estágio em virtude de ..., que será detalhado melhor as subseções seguintes.

\subsection{Foco do estudo do estágio}

A empresa possui um sistema de ...(ou não possui ou deveria possuir). Nele, os colaboradores da empresa podem .... Em razão disto, frequentemente os colaboradoes e/ou funcionários ... paravam de executar suas funções para verificar se existia ... o que poderia ser prejudicial para a empresa. Portanto, para solucionar este problema, foi proposto ... capaz de ... que proverá... 

	Dentro deste contexto, o principal foco do estágio foi o desenvolvimento de um estudo/análise/verificação/..., o qual ...utilizando o banco de dados que fica alocado no servidor local da empresa, para verificar se existia um novo chamado.
    
\begin{figure}[ht]
\centering
\includegraphics[width=.5\textwidth]{image01.jpg}
\caption{Figura do esquema proposto}
\label{fig:exampleFfig1}
\end{figure}

\section{A Intervenção proposta é de ...} \label{sec:intervencao}

O processo de Análise de Teste é de vital importância no desenvolvimento de um sistema, pois ele tem o objetivo de encontrar e depois corrigir erros no mesmo.

\subsection{Teste e Validação da proposta}
  
    Com o objetivo de verificar a eficiência e a eficácia da proposta de intervenção será necessário aplicar as seguintes técnicas/procedimentos...
   

\section{Considerações Finais} \label{sec:consideracoes}
Estudo, aprendizado e amadurecimento são três palavras que basicamente definem este estágio. Acredita-se que este foi de vital importância, pois contribuiu tanto para a minha vida profissional quando pessoal. Foi possível aplicar as teorias ensinadas em aula tanto na área de software quando na área de hardware, adquirindo uma ampla experiência e conhecimento na área da informática, além de conviver dentro de uma empresa, em um ambiente de trabalho onde tudo deve ser levado a sério, com responsabilidade e competência. 

	A empresa ... possui grande participação na região onde atual, principalmente ..., pois oferece aos seus clientes ... . Acredita-se ainda, que com as sugestões de melhoramento propostas nesse relatório possa ....


\section{Referencial Bibliográfico}\label{sec:referencial}

% estilo da bibliografia
\bibliographystyle{abbrv}
% chamando o arquivo refs.bib
\bibliography{refs}
\cite{horstmann2002core}


\end{document}
